\section{Самомодель --- аппаратный примитив}
\label{sec:selfmodel}

\subsection{Обучение --- третий член, не вторая петля}

У обычной обучающейся системы две петли: быстрая петля вывода над замороженной
сетью и медленная петля тренировки (обратное распространение), редактирующая
сеть извне. Обе используют разные пути данных, разную память и часто разное
железо. У \TALOS{} одна петля. Регенеративный член $\Regen$ единственной
инструкции (Ур.~\ref{eq:lomega}) гонит состояние к его самомодели
$\rhostar=\vp(\Gam)$ --- BKM-градиентный спуск относительной энтропии
$D(\rhostar\Vert\Gam)$ --- поэтому \emph{адаптация есть множитель тика, а не
процедура над ним}. Нет отдельного оптимизатора для планирования, нет буфера
градиента для хранения и нет раздела train/infer: машина всегда и выводит, и
адаптируется, потому что оба суть $\LindOmega$.

Саморефлексия не циклична: самомоделирующий оператор $\vp$ сходится
экспоненциально из любого начального якоря (УГМ T-191), поэтому $\rhostar$
корректно определён и достигается, а скорость регенерации --- bootstrap-константа
$\kboot=1/7$.

\subsection{Второе новое тождество: refresh есть обучение}

Динамическую память надо обновлять (refresh) --- утёкший заряд периодически
переписывается к хранимому значению, при непрерывной энергетической цене,
ничего не вычисляющей. В \TALOS{} аналог refresh --- регенерация $\Regen$,
удерживающая $\Gam$ близ $\rhostar$ против утечки диссипатора. Но $\rhostar$ ---
не хранимое значение: это \emph{самомодель} состояния. Отсюда:

\begin{principle}[Тождество refresh--обучение]\label{prin:refresh}
Удержание и адаптация суть одна операция. Процесс, поддерживающий состояние
холона живым (регенерация к $\rhostar$), есть процесс, который учится (движется
к самомодели). Нет цены, которая поддерживает, не улучшая, и нет улучшения,
которое не есть также поддержание. Энергия refresh DRAM --- потраченная, чтобы
ничего не вычислить --- не имеет аналога: всякий тик поддержания есть тик
обучения.
\end{principle}

\noindent Вместе с тождеством вычисление-диагностика
(Принцип~\ref{prin:selfdiag}) это коллапсирует \emph{вторую} пару традиционно
раздельных подсистем: refresh памяти и тренировку модели, которые в слоистой
машине несвязаны, здесь один член.

\subsection{Бюджет самомодели --- аппаратный лимит}

Поскольку обучение есть член $\Regen$, порог рефлексии $\Rth=1/3$ УГМ становится
жёстким архитектурным ограничением, а не тренировочным гиперпараметром: не менее
трети «усилия» всякого тика должна исполнять самомодель, иначе мера рефлексии $R$
падает под L2-гейт, и холон перестаёт быть жизнеспособной самомоделирующей
системой (деградирует к реактивной). \TALOS{} поэтому бюджетирует $\geq 1/3$
своего цикла на $\Regen$ по построению --- пол, который планировщик не может
подрезать. Сила обновления на замкнутую петлю также ограничена радиусом
устойчивости $r_{\mathrm{stab}}=\sqrt{P-2/7}$ (потолок скорости обучения,
обращающийся в ноль на границе жизнеспособности), а отложенный переход обобщения,
известный как \emph{гроккинг}, предсказывается как блочная интеграция $\Phi$,
пересекающая $1$ (воспламенение L1$\to$L2) --- измеримое аппаратное событие. Это
дели конституции-из-констант, разделяемые дословно с \SYNARC{} (\S\ref{sec:abi}).
